%! Complex Numbers — Unit 3, Lesson 3.4 — Homework (Answer Key)
\documentclass[10pt]{article}
\usepackage{algebra2-article}
\usepackage{algebra2-key}   % pulls in -boxes; do NOT also load -boxes
\newcommand{\TallMath}[1]{$\displaystyle #1\rule[-1.4em]{0pt}{3.2em}$}

\begin{document}

\pageheader{Unit 3, Lesson 3.4}{Homework --- Answer Key}
\namedateperiod

\vspace{0.10in}

\begin{notesbox}{Practice}
Show your work. Pull the $i$ out \emph{first}, replace $i^2=-1$ when you multiply, and write every answer
in standard $a+bi$ form.

\begin{enumerate}[leftmargin=1.9em, itemsep=10pt]
  \item \textbf{Rewrite with $i$} (simplest radical form).
    \[ \sqrt{-25}=\ans{$5i$} \quad \sqrt{-8}=\ans{$2i\sqrt{2}$} \quad \sqrt{-72}=\ans{$6i\sqrt{2}$} \quad \sqrt{-100}=\ans{$10i$} \quad -\sqrt{-49}=\ans{$-7i$} \]

  \item \textbf{Write in $a+bi$ form.}
    \[ 4+\sqrt{-9}=\ans{$4+3i$} \quad -2-\sqrt{-16}=\ans{$-2-4i$} \quad \sqrt{-121}=\ans{$0+11i$} \quad 10=\ans{$10+0i$} \]

  \item \textbf{Add or subtract.}
    \begin{enumerate}[label=(\alph*), leftmargin=1.9em, itemsep=5pt]
      \item $(5+3i)+(2-7i)=$ \ans{$7-4i$}
      \item $(-4+i)-(3-6i)=$ \ans{$-7+7i$}
      \item $(8-2i)-(8+2i)=$ \ans{$-4i$}
    \end{enumerate}

  \item \textbf{Multiply.} FOIL, then replace $i^2=-1$.
    \begin{enumerate}[label=(\alph*), leftmargin=1.9em, itemsep=5pt]
      \item $(2+i)(3+4i)=$ \ans{$2+11i$}
      \item $(5-3i)(2-i)=$ \ans{$7-11i$}
      \item $(1+6i)^2=$ \ans{$-35+12i$}
      \item $3i(4-5i)=$ \ans{$15+12i$}
    \end{enumerate}

  \item \textbf{Conjugate products.} Multiply by the conjugate; your answer should be a real number.
    \[ (7+2i)(7-2i)=\ans{$53$} \qquad (3-5i)(3+5i)=\ans{$34$} \]

  \item \textbf{Powers of $i$.} Reduce the exponent mod $4$.
    \[ i^{12}=\ans{$1$} \quad i^{27}=\ans{$-i$} \quad i^{50}=\ans{$-1$} \quad i^{101}=\ans{$i$} \]

  \item \textbf{Solve over the complex numbers.} Both solutions, in $a+bi$ form.
    \[ x^2=-64:\ x=\ans{$\pm8i$} \qquad x^2+20=0:\ x=\ans{$\pm2i\sqrt{5}$} \]

  \item \textbf{Explain.} A classmate claims ``$i^2=1$, because squaring always makes a number positive.''
        What is wrong with this reasoning, and what is $i^2$ really?
        \ansline{``Squaring gives a nonnegative result'' is a rule for \emph{real} numbers. $i$ is not real; it is \emph{defined} so that $i^2=-1$. So $i^2=-1$, not $1$.}
\end{enumerate}
\end{notesbox}

\begin{extensionbox}
\textbf{Complex roots by completing the square.} Solve each; the two solutions are complex conjugates.
\begin{enumerate}[label=(\alph*), leftmargin=1.8em, itemsep=5pt]
  \item $x^2-6x+13=0 \Rightarrow (x-3)^2=$ \ans{$-4$}, so $x=$ \ans{$3\pm2i$}
  \item $2x^2+8=0$ (divide by $2$ first) $\Rightarrow x=$ \ans{$\pm2i$}
  \item Check (a): substitute one of your roots back into $x^2-6x+13$ and show it gives $0$.
\end{enumerate}
\ansline{Using $x=3+2i$: $(3+2i)^2-6(3+2i)+13=(5+12i)-(18+12i)+13=0$. \checkmark}
\end{extensionbox}

\begin{spiralbox}
\textbf{Coming next --- Lesson 3.5: The Quadratic Formula \& the Discriminant.} You can now handle a
\emph{negative} under a square root, so \emph{every} quadratic is solvable. Next you'll get one formula
that solves them all, $x=\dfrac{-b\pm\sqrt{b^2-4ac}}{2a}$, and use the \textbf{discriminant} $b^2-4ac$ to
predict --- before solving --- whether the roots are two real, one real, or two complex conjugates.
\end{spiralbox}

\begin{teachernote}
Answer traps to flag while grading: \#1 not simplifying ($\sqrt{-72}=6i\sqrt2$, not $\sqrt{72}\,i$) or
mishandling the leading minus in $-\sqrt{-49}=-7i$; \#4c forgetting $i^2=-1$ so $(1+6i)^2$ stops at
$1+12i+36i^2$; \#6 forgetting the remainder ($i^{12}=1$ because $12$ is a multiple of $4$); \#7 dropping
the $\pm$. The extension is the bridge to 3.5 --- if a student asks \emph{why} the roots are conjugates,
that ``why'' (the discriminant $b^2-4ac<0$) is exactly next lesson.
\end{teachernote}

\end{document}
